% DMRN+21: the official 2026 kit ships the 2025 style; its event header and
% footer are updated in the local copy, without changing its page layout.
\documentclass[a4paper]{article}
\usepackage[T1]{fontenc}
\usepackage{caption}
\usepackage{times,amsmath,graphicx,url,dmrn2026}
\usepackage[table]{xcolor}
\definecolor{tabpurple}{HTML}{7625A8}
\usepackage{indentfirst}
\setlength{\parskip}{0.1em}
\captionsetup[figure]{font=footnotesize,labelfont=footnotesize}
\captionsetup[table]{font=footnotesize,labelfont=footnotesize}
\setlength{\belowcaptionskip}{0pt}
\renewcommand{\thesection}{\Roman{section}}

\title{Fast Backpropagation for a Finite-Difference Time-Domain Kirchhoff--Carrier String}
\name{Pablo Tablas de Paula, Emmanouil D. Benetos and Joshua D. Reiss\thanks{This work was supported by UKRI grant EP/S022694.}}
\address{Centre for Digital Music, Queen Mary University of London, UK}

\begin{document}
\tenpt
\maketitle
\begin{sloppy}
\begin{abstract}
We accelerate backpropagation through a nonlinear finite-difference string
using CUDA state reconstruction, checkpoints and three scalars per step.
Russo et al.'s time step is preserved; a fourth-order spatial stencil removes
the grid's dispersion.
\end{abstract}

\section{Introduction}
Plucked-string synthesis spans Karplus--Strong, digital waveguides, modal and
FDTD methods; FDTD captures the string's nonlinearity and its fret and finger
collisions, and energy-based schemes now run full guitar models in real
time~\cite{bilbao2024guitar}. Differentiable implementations learn a model's
parameters from recordings by gradient descent~\cite{engel2020ddsp,tablas2026sound}.
Backpropagating through an FDTD is costly because reverse mode keeps every
node at every step, so differentiable physical models have favoured modal
formulations with exact linear recurrences~\cite{diaz2025}. We show that an
FDTD string can be backpropagated at modal cost.

\section{Model and method}
We use Russo et al.'s non-SAV reference scheme~\cite{russo2026}. For
displacement $u(x,t)$ and $\mu=\rho A$,
\begin{equation}
\begin{aligned}
\mu u_{tt}={}&T_0u_{xx}-EIu_{xxxx}-2\mu\sigma_0u_t\\
 &+\frac{EA}{2L}\!\left(\int_0^L u_x^2\,dx\right)u_{xx}
 +2\mu\sigma_1u_{txx},
\end{aligned}
\label{eq:pde}
\end{equation}
with length $L$, density $\rho$, area $A$, modulus $E$, second moment $I$,
tension $T_0$ and damping $\sigma_0,\sigma_1$.

The reference step uses a Sherman--Morrison rank-one solve; a full sine
transform diagonalizes the spatial operators without dropping modes. With
$p,u,v$ successive states, $\ell=Bu+Pp$, $v=m(\ell-q\alpha)$, $B,P$ diagonal,
$m=[\mu(1+\sigma_0k)]^{-1}$, $q=(k/2)\sqrt{hEA/L}\,D_{xx}u$ for timestep $k$
and spacing $h$, and $\alpha=(q^Tp+mq^T\ell)/(1+mq^Tq)$, the previous state
is $p=P^{-1}(v/m-Bu+q\alpha)$. The forward pass stores both states every 256
steps plus $\alpha$, $q^Tq$ and $q^T\ell$ per step; the reverse pass
reconstructs each segment and differentiates all dependencies, CUDA-parallel
over modes and batch items.

The reference's three-point $D_{xx}$ disperses, drifting the waveform 11.7\%
from the converged solution within a second; we use the five-point stencil
$(-1,16,-30,16,-1)/12h^2$, closed at the simply supported ends by odd
reflection, with $D_{xxxx}$ its square. It is diagonal in the same sine
basis, so only the eigenvalues change, $-4s^2/h^2\to-4s^2(1+s^2/3)/h^2$ with
$s=\sin(\pi n/2N)$; time step, rank-one solve and adjoint are untouched, and
the stability bound is the reference's with $h\to h\sqrt{3}/2$.

\section{Evaluation}
We compare our cached adjoint (\emph{Ours}), \emph{Autodiff} (direct PyTorch
autodiff) and \texttt{jaxdiffmodal}~\cite{diaz2025} on an A100 80 GB:
1 s, float64, identical pickup, loss and the reference code's raised-cosine
pluck; FDTD on 176 intervals at 88.2 kHz (2$\times$, decimated), within
1.2\% waveform relative $L^2$ of the modal solution, and \texttt{jaxdiffmodal}
with 100 continuous modes at 44.1 kHz.
In Fig.~\ref{fig:benchmark} our time is nearly flat in batch size and our
memory grows slowly; \texttt{jaxdiffmodal} is slightly slower and its memory
grows linearly until it runs out; Autodiff is nearly three orders of magnitude
slower and runs out of memory first.
The three- and five-point variants match the (correspondingly modified)
reference code to $2.5\times10^{-9}$ and $5.0\times10^{-10}$ relative $L^2$
over 4 s, and all eight gradients pass a 2 ms finite-difference check.

\begin{figure}[!h]
 \centering
 \includegraphics[width=\columnwidth]{figures/benchmark-dots-label.pdf}
 \caption{Raised-cosine pluck, 1 s, eight gradients; numbers: seconds per evaluation.}
 \label{fig:benchmark}
\end{figure}

\bibliographystyle{IEEEtran}
{\footnotesize\bibliography{references}}
\end{sloppy}
\end{document}
